\documentclass[11pt]{article}

\usepackage[T1]{fontenc}
\usepackage[utf8]{inputenc}
\usepackage{amsmath,amssymb,amsthm}
\usepackage{booktabs}
\usepackage[margin=1in]{geometry}
\usepackage[hidelinks]{hyperref}
\hypersetup{
  pdftitle={Schur-log-concavity for the complete eight-vertex layer of the multiweighted-bond symmetric function},
  pdfauthor={Xiaojian Zeng},
  pdfsubject={Exact verification of Schur-log-concavity for connected graphs on eight vertices},
  pdfkeywords={symmetric functions, Schur positivity, graph invariants, exact computation}
}

\newtheorem{theorem}{Theorem}
\newtheorem{lemma}[theorem]{Lemma}
\newcommand{\Sym}{\mathrm{Sym}}
\newcommand{\coeff}[2]{[#1]#2}

\title{Schur-log-concavity for the complete eight-vertex layer\\
of the multiweighted-bond symmetric function}
\author{Xiaojian Zeng\\
\small PBC School of Finance, Tsinghua University, Beijing 100190, China\\
\small \texttt{xiaojianzeng@foxmail.com}}
\date{6 September 2026}

\begin{document}
\maketitle

\begin{abstract}
Gonz\'alez D'Le\'on and Wachs introduced a symmetric-function analogue
$\Psi_G$ of the chromatic polynomial and conjectured that it is
Schur-log-concave.  They reported exhaustive verification for connected
graphs on at most seven vertices.  We settle the next complete finite layer:
for every connected simple graph $G$ on eight vertices and every $d\geq 1$,
the difference
\[
  (\Psi_G|_{\deg d})^2
  -(\Psi_G|_{\deg(d-1)})(\Psi_G|_{\deg(d+1)})
\]
is Schur-positive.  An exact computation covers all $11{,}117$ isomorphism
classes and all $294$ Schur coefficients arising in the seven nonzero
degrees for each class.  In fact every one of the $3{,}268{,}398$
coefficients tested is strictly positive, with global minimum $1$; all
higher-degree differences vanish.  The computation uses an exhaustive
atlas-extension argument, a connected-partition recurrence, and integer
arithmetic in the complete-homogeneous and Schur bases.
\end{abstract}

\section{Statement and context}

Let $G=(V,E)$ be a finite simple graph.  A \emph{bond partition} of $G$ is
a set partition of $V$ whose blocks induce connected subgraphs.  Write
$\Pi_G$ for the set of bond partitions.  We use the multiweighted-bond
symmetric function $\Psi_G\in\Sym$ and the M\"obius symmetric function $M_H$
of Gonz\'alez D'Le\'on and Wachs~\cite{gonzalez-wachs2026}.  Put
\[
  F_{G,d}=\Psi_G|_{\deg d},\qquad
  L_{G,d}=F_{G,d}^2-F_{G,d-1}F_{G,d+1}.
\]
The Schur-log-concavity conjecture in~\cite[Conjecture~7.3]{gonzalez-wachs2026}
asserts that $L_{G,d}$ is Schur-positive for every graph and every $d$.
It refines the ordinary chromatic-polynomial log-concavity proved by
Huh~\cite{huh2012}.

\begin{theorem}\label{thm:main}
For every connected simple graph $G$ on exactly eight vertices and every
integer $d\geq1$, the symmetric function $L_{G,d}$ is Schur-positive.
More precisely, for $1\leq d\leq7$ every coefficient in
\[
  L_{G,d}=\sum_{\lambda\vdash 2d}c_{G,d,\lambda}s_\lambda
\]
is a positive integer, and the minimum of all these coefficients over all
$G,d,\lambda$ is $1$.  For $d\geq8$, $L_{G,d}=0$.
\end{theorem}

This result concerns only the complete eight-vertex layer; it does not settle
the conjecture for arbitrary order.  The source proves alternating
$e$-positivity of the individual homogeneous pieces for chordal graphs, but
that result does not imply Schur positivity of the quadratic differences.
Accordingly, chordal and nonchordal graphs were both computed.  Table
\ref{tab:minima} gives the number of classes in each stratum and the smallest
Schur coefficient at each degree.

\begin{table}[ht]
\centering
\caption{Exact coefficient minima by stratum.}
\label{tab:minima}
\begin{tabular}{lrrccccccc}
\toprule
 & classes & coefficients & $d=1$ & $2$ & $3$ & $4$ & $5$ & $6$ & $7$ \\
\midrule
chordal    & $1{,}614$ & $474{,}516$   & $7$  & $158$ & $490$  & $490$  & $196$ & $28$  & $1$ \\
nonchordal & $9{,}503$ & $2{,}793{,}882$ & $18$ & $288$ & $1176$ & $1695$ & $946$ & $186$ & $9$ \\
\bottomrule
\end{tabular}
\end{table}

The closest prior result is the source computation through seven vertices.
A search for the exact conjecture and the follow-up announced in that source,
performed through 6 September 2026, located no separate primary paper settling
the eight-vertex layer.  This is a dated, bounded literature comparison, not a
claim of priority.  For nearby computational context, Zeng~\cite{zeng2026}
uses selected coefficients of the classical chromatic symmetric function and
an exact cut-partition implementation to reconstruct a class of trees.  That
work concerns a different invariant and provides no evidence for
Schur-log-concavity here.

\section{Exact connected-partition recurrence}

For $S\subseteq V$ and $k\geq0$, define
\[
 A_G(S,k)=
 \sum_{\substack{\pi\in\Pi_{G[S]}\\|\pi|=k}}
 \prod_{B\in\pi}M_{G[B]},
 \qquad A_G(\varnothing,0)=1.
\]
All other impossible cases are zero.  Formula (7.1) and Theorem 7.2 of
\cite{gonzalez-wachs2026} give, for $n=|V|$,
\begin{equation}\label{eq:FfromA}
  F_{G,d}=A_G(V,n-d).
\end{equation}

\begin{lemma}\label{lem:recurrence}
Fix an order on $V$, and let $a$ be the least vertex of nonempty $S$.  For
$k\geq2$,
\begin{equation}\label{eq:anchored}
 A_G(S,k)=
 \sum_{\substack{a\in B\subsetneq S\\G[B]\text{ connected}}}
 M_{G[B]}A_G(S\setminus B,k-1).
\end{equation}
If $G[S]$ is connected and $|S|\geq2$, then
\begin{equation}\label{eq:mobius}
 M_{G[S]}=-\sum_{k=2}^{|S|}h_{k-1}A_G(S,k).
\end{equation}
Together with $M_{G[\{v\}]}=1$, these equations compute every term in
\eqref{eq:FfromA} by increasing $|S|$.
\end{lemma}

\begin{proof}
In a connected-block partition of $S$, exactly one block contains $a$.
Removing it gives a connected-block partition of its complement, and this is
a bijection proving \eqref{eq:anchored}.  Theorem 5.7, formula (5.3), of
\cite{gonzalez-wachs2026} expresses $M_{G[S]}$ as the negative sum over all
nontrivial bond partitions, with the factor $h_{k-1}$ for a partition with
$k$ blocks.  Grouping that formula by $k$ yields \eqref{eq:mobius}.  Every
block in \eqref{eq:anchored} is proper, so increasing subset cardinality is a
valid induction.
\end{proof}

The implementation represents a symmetric function as a sparse integer
linear combination of $h_\mu=\prod_i h_{\mu_i}$.  Products in this basis
concatenate and sort the parts.  Conversion to the Schur basis uses
\begin{equation}\label{eq:kostka}
  h_\mu=\sum_{\lambda\vdash |\mu|}K_{\lambda\mu}s_\lambda,
\end{equation}
where the nonnegative integers $K_{\lambda\mu}$ are generated by successive
applications of the horizontal-strip Pieri rule; see
Macdonald~\cite[Chapter~I]{macdonald1995}.  Thus the recurrence, products,
basis conversion, and sign tests use integer arithmetic only.

As a definition-level regression test, direct enumeration of the finite
two-variable multiweighted-bond poset for the path $P_3$ agrees exactly with
Lemma~\ref{lem:recurrence}.  It gives
\[
 \Psi_{P_3}=1-2e_1+e_{1,1}+e_2,
 \qquad L_{P_3,1}=3s_2+2s_{1,1}.
\]
The last two coefficients are interchanged in the displayed equality in
\cite[Example~7.4]{gonzalez-wachs2026}; the preceding expression for
$\Psi_{P_3}$ in that paper gives the value above.  This example-level swap
does not affect positivity or the eight-vertex calculation.

\section{Exhaustive graph enumeration}

The graph atlas of Read and Wilson~\cite{read-wilson1998}, exposed by
NetworkX~\cite{hagberg-schult-swart2008}, contains one representative of every
isomorphism class on seven vertices; there are $1{,}044$ such representatives.
For every representative $H$, the program adjoins a new vertex and chooses
each of the $2^7-1=127$ nonempty subsets of $V(H)$ as its neighborhood.  After
discarding disconnected extensions, $119{,}980$ labeled candidates remain.

Candidates are first bucketed by the degree multiset, triangle-count
multiset, and a four-round Weisfeiler--Lehman graph hash.  These are used only
as isomorphism invariants.  Inside each bucket, equality is decided by
NetworkX's exact graph-isomorphism routine, based on the VF2 algorithm
\cite{cordella-et-al2004}; a hash is never used to declare two graphs
isomorphic.  Exactly $11{,}117$ representatives remain.

\begin{lemma}\label{lem:graphs}
The retained family contains exactly one representative of every isomorphism
class of connected simple graphs on eight vertices.
\end{lemma}

\begin{proof}
Let $G$ be such a graph and choose a vertex $v$.  The seven-vertex graph
$G-v$, possibly disconnected, is isomorphic to an atlas representative $H$.
Since $G$ is connected, $N_G(v)$ is nonempty.  Transporting $N_G(v)$ through
the isomorphism from $G-v$ to $H$ shows that an extension isomorphic to $G$
is generated.  Isomorphic candidates have equal bucket invariants and are
therefore compared by the exact test.  Conversely, the exact test removes a
candidate only when it is isomorphic to one already retained.  This proves
existence and uniqueness.
\end{proof}

\section{Certificate and proof of the theorem}

For every retained graph and every $d=1,\ldots,7$, the program computes
$F_{G,d-1},F_{G,d},F_{G,d+1}$ from Lemma~\ref{lem:recurrence}, forms
$L_{G,d}$ in the $h$ basis, and converts it using \eqref{eq:kostka}.  It then
iterates over \emph{every} partition of $2d$, inserting zero for a coefficient
absent from the sparse representation before taking the minimum.  The numbers
of tested coordinates are
\[
 p(2)+p(4)+p(6)+p(8)+p(10)+p(12)+p(14)=294
\]
per graph and $11{,}117\cdot294=3{,}268{,}398$ in total.  The resulting
minimum values are those in Table~\ref{tab:minima}; all are positive.

Lemma~\ref{lem:graphs} supplies every graph class.  Lemma
\ref{lem:recurrence} supplies the exact homogeneous pieces, and the Pieri
conversion supplies their complete Schur expansions.  Hence the positive
minimum certifies the assertion for $1\leq d\leq7$.  Because a connected
eight-vertex graph has $F_{G,d}=0$ outside $0\leq d\leq7$, both terms in
$L_{G,d}$ vanish for $d\geq8$.  This proves Theorem~\ref{thm:main}.

The supplementary record contains one graph6 representative per line, its
chordality flag, and the degree-by-degree coefficient summaries.  A clean
replay under Python 3.10.12 and NetworkX 3.4.2 regenerated all records
byte-for-byte.  From the artifact root, the full calculation is reproduced by
\begin{verbatim}
python3 evidence/eight_vertex_scan.py \
  --records /tmp/eight_vertex_records.replay.jsonl \
  --summary /tmp/eight_vertex_summary.replay.json \
  --progress-every 0
\end{verbatim}
A complete run required only integer arithmetic and finished in about 136
seconds in the documented environment.  These finite data do not suggest a
proof for graphs on nine or more vertices, and strict positivity in this layer
must not be extrapolated beyond the stated range.

\paragraph{Computational and AI assistance.}
Python and NetworkX were used for graph enumeration, exact isomorphism tests,
and the integer symmetric-function computation.  OpenAI Codex assisted with
implementation, proof organization, manuscript preparation, and a separate
audit that included a full deterministic replay.  Model output was not treated
as mathematical evidence: the theorem rests on the explicit recurrence,
exhaustive enumeration argument, exact code, certificate, and replay described
above.

\bibliographystyle{plain}
\bibliography{references}

\end{document}
